% Arsitektur sistem penuh Rupiah Digital: client → backend → Fabric → peers
% Di-input dari BAB III §3.1.3

\begin{figure}[H]
\centering
\resizebox{\textwidth}{!}{%
\begin{tikzpicture}[
  appbox/.style={draw=black, fill=white, rectangle, rounded corners=3pt,
    minimum width=5.0cm, minimum height=1.1cm, align=center, font=\small},
  gwbox/.style={draw=black, fill=black!5, rectangle, rounded corners=3pt,
    minimum width=5.0cm, minimum height=1.1cm, align=center, font=\small},
  pgbox/.style={draw=black, fill=blue!8, rectangle, rounded corners=3pt,
    minimum width=3.8cm, minimum height=1.1cm, align=center, font=\small},
  peer/.style={draw=black, fill=white, rectangle, rounded corners=2pt,
    minimum width=2.9cm, minimum height=1.1cm, align=center, font=\scriptsize},
  obsbox/.style={draw=black, dashed, fill=white, rectangle, rounded corners=2pt,
    minimum width=2.9cm, minimum height=1.1cm, align=center, font=\scriptsize},
  ordnode/.style={draw=black, fill=black!8, rectangle, rounded corners=3pt,
    minimum width=3.8cm, minimum height=1.1cm, align=center, font=\scriptsize\bfseries},
  cdb/.style={draw=black, fill=gray!20, rectangle,
    minimum width=2.5cm, minimum height=0.75cm, align=center, font=\scriptsize},
  cabox/.style={draw=black, fill=yellow!20, rectangle, rounded corners=2pt,
    minimum width=2.5cm, minimum height=0.75cm, align=center, font=\tiny},
  ar/.style={-{Stealth}, thick},
  bar/.style={{Stealth}-{Stealth}, thick},
  lbl/.style={font=\scriptsize, midway},
]

%── Lapisan Aplikasi ──────────────────────────────────────────────────────────

\node[appbox] (fe) at (6.0, 16.5) {\textbf{Client}\\Web UI / Client App};
\node[appbox] (be) at (6.0, 14.5) {\textbf{Backend API}\\Go REST + Gorilla Mux + JWT/RBAC};
\node[pgbox]  (pg) at (0.0, 14.5) {\textbf{PostgreSQL}\\KYC \textit{off-chain}};
\node[gwbox]  (gw) at (6.0, 12.5) {\textbf{Fabric Gateway SDK v1.8}\\gRPC / TLS};

\draw[ar]  (fe.south) -- (be.north)  node[lbl, right]{HTTP/REST};
\draw[bar] (be.west)  -- (pg.east)   node[lbl, above]{SQL};
\draw[ar]  (be.south) -- (gw.north)  node[lbl, right]{Gateway API};

%── Jaringan Hyperledger Fabric ───────────────────────────────────────────────
% Latar belakang kotak digambar PERTAMA agar node internal tampak di atasnya

\fill[rounded corners=8pt, gray!5]  (-2.8, 2.2) rectangle (14.8, 11.0);
\draw[rounded corners=8pt, black]   (-2.8, 2.2) rectangle (14.8, 11.0);
\node[font=\footnotesize\bfseries, anchor=south] at (6.0, 11.0)
  {\textbf{Hyperledger Fabric 2.5}\quad
   (digital-rupiah channel, kebijakan dukungan mayoritas, TLS)};

% Panah dari Gateway ke batas atas kotak jaringan
\draw[ar] (gw.south) -- (6.0, 11.0)
  node[lbl, right]{\textit{endorse} + submit (gRPC)};

%── Orderer ───────────────────────────────────────────────────────────────────

\node[ordnode] (ord) at (6.0, 9.5)
  {orderer.paynet:7050\\Raft CFT};

%── Lima Peer ─────────────────────────────────────────────────────────────────

\node[peer]   (pbi)  at (0.0,  7.0)
  {peer0.bi\\BankIndonesiaOrg\\:7051};
\node[peer]   (phim) at (3.0,  7.0)
  {peer0.himbara\\HimbaraBankOrg\\:9051};
\node[peer]   (pcom) at (6.0,  7.0)
  {peer0.commercial\\CommercialBankOrg\\:11051};
\node[peer]   (ppjp) at (9.0,  7.0)
  {peer0.pjp\\PJPOrg\\:13051};
\node[obsbox] (pojk) at (12.0, 7.0)
  {peer0.ojk\\OJKObserverOrg\\:12051};

% Orderer → semua peer (pengiriman blok)
\draw[ar] (ord.south) to[out=-110,in=80] (pbi.north);
\draw[ar] (ord.south) to[out=-100,in=80] (phim.north);
\draw[ar] (ord.south) -- (pcom.north);
\draw[ar] (ord.south) to[out=-80,in=100]  (ppjp.north);
\draw[ar] (ord.south) to[out=-70,in=100]  (pojk.north);

% peer0.bi → orderer (endorsement submit)
\draw[ar, dashed] (pbi.north) to[out=70,in=-120]
  node[near end, left, font=\tiny]{submit} (ord.west);

%── CouchDB ───────────────────────────────────────────────────────────────────

\node[cdb] (cbi)  at (0.0,  5.5) {CouchDB :5984};
\node[cdb] (chim) at (3.0,  5.5) {CouchDB :6984};
\node[cdb] (ccom) at (6.0,  5.5) {CouchDB :7984};
\node[cdb] (cpjp) at (9.0,  5.5) {CouchDB :9984};
\node[cdb] (cojk) at (12.0, 5.5) {CouchDB :8984};

\draw[ar] (pbi)  -- (cbi);
\draw[ar] (phim) -- (chim);
\draw[ar] (pcom) -- (ccom);
\draw[ar] (ppjp) -- (cpjp);
\draw[ar] (pojk) -- (cojk);

%── CA per Organisasi ─────────────────────────────────────────────────────────

\node[cabox] (cabi)  at (0.0,  4.0) {CA\\BankIndonesia};
\node[cabox] (cahim) at (3.0,  4.0) {CA\\HimbaraBank};
\node[cabox] (cacom) at (6.0,  4.0) {CA\\CommercialBank};
\node[cabox] (capjp) at (9.0,  4.0) {CA\\PJP};
\node[cabox] (caojk) at (12.0, 4.0) {CA\\OJKObserver};

\draw[bar] (cabi)  -- (cbi);
\draw[bar] (cahim) -- (chim);
\draw[bar] (cacom) -- (ccom);
\draw[bar] (capjp) -- (cpjp);
\draw[bar] (caojk) -- (cojk);

%── Orderer CA (di bawah kotak jaringan) ─────────────────────────────────────

\node[cabox, minimum width=4.0cm] (caord) at (6.0, 1.0)
  {BankIndonesia Orderer CA};
\draw[ar] (6.0, 2.2) -- (caord.north);

\end{tikzpicture}%
}
\caption{Arsitektur sistem Rupiah Digital: lapisan \textit{client}, \textit{backend API} (Go REST + JWT/RBAC), basis data PostgreSQL luring, \textit{Fabric Gateway SDK}, serta jaringan Hyperledger Fabric 2.5 dengan lima organisasi \textit{peer}, CouchDB, CA per organisasi, dan \textit{orderer} Raft}
\label{fig:arsitektur}
\end{figure}
